% GENERATED FILE. Edit canonical lessons, then run npm run generate:books.
% canonical-source-hash: fnv1a64:b236c1b39e1ad1f3
% canonical-lessons: UR-C510-chhilna-peel, UR-C510-kutna, UR-C510-okhli, UR-C510-pisvana, UR-C510-nichorna

\chapter{Peel, Pound, Squeeze}
\label{ch:ur-peel-pound-squeeze}
\hlchaptermodality{510}

\begin{chapteropening}
\textbf{By the end of this chapter:} \emph{I can describe getting food ready to cook.}

Complete the last lesson of chapter 510: I can describe getting food ready to cook.
\end{chapteropening}

\section[\texorpdfstring{\ur{چھیلنا} (chhīlnā)}{چھیلنا (chhīlnā)}]{\ur{چھیلنا} (chhīlnā) — to peel}

\label{lesson:UR-C510-chhilna-peel}



\begin{quote}
Before the new one: say the Urdu for an earthen cooking pot, then the Urdu for to spill.
\end{quote}

\subsection*{You'll want to know: \ur{چھیلنا}}
\textbf{\ur{چھیلنا}} — \emph{chhīlnā} — \textquotedblleft{}to peel\textquotedblright{}.

\begin{quote}
\ur{میں} \ur{نے} \ur{چھری} \ur{سے} \ur{آلو} \ur{چھیلے۔}
\end{quote}

\begin{quote}
\emph{maiṅ ne chhurī se ālū chhīle.}
\end{quote}

\begin{quote}
I peeled the potatoes with a knife.
\end{quote}

Past, with \ur{نے} (he or she …): \ur{اس} \ur{نے} \ur{سیب} \ur{چھیلا}, \emph{us ne seb chhīlā}.

\begin{cousinweb}
\textbf{\ur{چھیلنا}} (\emph{chhīlnā}) is inherited: CDIAL lists the Hindi cousin \emph{chīlnā}, \textquotedblleft{}to peel\textquotedblright{}, under \emph{*chilla}, a form it reconstructs (CDIAL 5051).
\end{cousinweb}

\begin{grammarlens}[title={What you've built}]
One more word for describing what you lived and felt.
\end{grammarlens}

\subsection*{Your turn}
\begin{itemize}
\raggedright
  \item \emph{Say it:} \emph{chhīlnā}
  \item \emph{Say it:} \emph{chhīlnā}, once more
  \item \emph{Say it:} \emph{maiṅ ne chhurī se ālū chhīle.} — I peeled the potatoes with a knife.
  \item \emph{Recall:} say the Urdu for an earthen cooking pot, then the Urdu for to spill, then say \emph{chhīlnā} again
\end{itemize}

\begin{tcolorbox}[breakable,colback=teal!4,colframe=teal!35!black,title={Before you move on}]
What does \ur{چھیلنا} mean? (\textquotedblleft{}To peel\textquotedblright{}.) Say it once more.
\end{tcolorbox}
\section[\texorpdfstring{\ur{کوٹنا} (kūṭnā)}{کوٹنا (kūṭnā)}]{\ur{کوٹنا} (kūṭnā) — to pound}

\label{lesson:UR-C510-kutna}



\begin{quote}
Before the new one: say the Urdu for to spill, then the Urdu for to peel.
\end{quote}

\subsection*{You'll want to know: \ur{کوٹنا}}
\textbf{\ur{کوٹنا}} — \emph{kūṭnā} — \textquotedblleft{}to pound\textquotedblright{}.

\begin{quote}
\ur{دادی} \ur{نے} \ur{کالی} \ur{مرچ} \ur{کوٹی۔}
\end{quote}

\begin{quote}
\emph{dādī ne kālī mirch kūṭī.}
\end{quote}

\begin{quote}
Grandmother pounded the black pepper.
\end{quote}

Past, with \ur{نے} (he or she …): \ur{اس} \ur{نے} \ur{مسالہ} \ur{کوٹا}, \emph{us ne masālā kūṭā}.

\begin{cousinweb}
\textbf{\ur{کوٹنا}} (\emph{kūṭnā}) is inherited from Sanskrit \emph{kuṭṭayati}, \textquotedblleft{}pounds, threshes\textquotedblright{} (CDIAL 3241).
\end{cousinweb}

\begin{grammarlens}[title={What you've built}]
One more word for describing what you lived and felt.
\end{grammarlens}

\subsection*{Your turn}
\begin{itemize}
\raggedright
  \item \emph{Say it:} \emph{kūṭnā}
  \item \emph{Say it:} \emph{kūṭnā}, once more
  \item \emph{Say it:} \emph{dādī ne kālī mirch kūṭī.} — Grandmother pounded the black pepper.
  \item \emph{Recall:} say the Urdu for to spill, then the Urdu for to peel, then say \emph{kūṭnā} again
\end{itemize}

\begin{tcolorbox}[breakable,colback=teal!4,colframe=teal!35!black,title={Before you move on}]
What does \ur{کوٹنا} mean? (\textquotedblleft{}To pound\textquotedblright{}.) Say it once more.
\end{tcolorbox}
\section[\texorpdfstring{\ur{اوکھلی} (okhlī)}{اوکھلی (okhlī)}]{\ur{اوکھلی} (okhlī) — a mortar (for pounding)}

\label{lesson:UR-C510-okhli}



\begin{quote}
Before the new one: say the Urdu for to peel, then the Urdu for to pound.
\end{quote}

\subsection*{You'll want to know: \ur{اوکھلی}}
\textbf{\ur{اوکھلی}} — \emph{okhlī} — \textquotedblleft{}a mortar (for pounding)\textquotedblright{}.

\begin{quote}
\ur{دادی} \ur{اوکھلی} \ur{میں} \ur{مسالہ} \ur{کوٹتی} \ur{تھیں۔}
\end{quote}

\begin{quote}
\emph{dādī okhlī meṅ masālā kūṭtī thīṅ.}
\end{quote}

\begin{quote}
Grandmother used to pound spices in a mortar.
\end{quote}

Feminine.

\begin{cousinweb}
\textbf{\ur{اوکھلی}} (\emph{okhlī}) is inherited: CDIAL lists the Hindi cousin \emph{okhlī}, \textquotedblleft{}a small wooden mortar\textquotedblright{}, under Sanskrit \emph{ulūkhala}, \textquotedblleft{}a mortar\textquotedblright{} (CDIAL 2360).
\end{cousinweb}

\begin{grammarlens}[title={What you've built}]
One more word for describing what you lived and felt.
\end{grammarlens}

\subsection*{Your turn}
\begin{itemize}
\raggedright
  \item \emph{Say it:} \emph{okhlī}
  \item \emph{Say it:} \emph{okhlī}, once more
  \item \emph{Say it:} \emph{dādī okhlī meṅ masālā kūṭtī thīṅ.} — Grandmother used to pound spices in a mortar.
  \item \emph{Recall:} say the Urdu for to peel, then the Urdu for to pound, then say \emph{okhlī} again
\end{itemize}

\begin{tcolorbox}[breakable,colback=teal!4,colframe=teal!35!black,title={Before you move on}]
What does \ur{اوکھلی} mean? (\textquotedblleft{}A mortar (for pounding)\textquotedblright{}.) Say it once more.
\end{tcolorbox}
\section[\texorpdfstring{\ur{پسوانا} (pisvānā)}{پسوانا (pisvānā)}]{\ur{پسوانا} (pisvānā) — to have (something) ground}

\label{lesson:UR-C510-pisvana}



\begin{quote}
Before the new one: say the Urdu for to pound, then the Urdu for a mortar (for pounding).
\end{quote}

\subsection*{You'll want to know: \ur{پسوانا}}
\textbf{\ur{پسوانا}} — \emph{pisvānā} — \textquotedblleft{}to have (something) ground\textquotedblright{}.

\begin{quote}
\ur{امی} \ur{نے} \ur{گندم} \ur{کا} \ur{آٹا} \ur{پسوایا۔}
\end{quote}

\begin{quote}
\emph{ammī ne gandum kā āṭā pisvāyā.}
\end{quote}

\begin{quote}
Mum had the wheat ground into flour.
\end{quote}

Past, with \ur{نے} (he or she …): \ur{اس} \ur{نے} \ur{مسالہ} \ur{پسوایا}, \emph{us ne masālā pisvāyā}.

\begin{cousinweb}
\textbf{\ur{پسوانا}} (\emph{pisvānā}) is the second causative of \emph{pīsnā}, \textquotedblleft{}to grind\textquotedblright{}, which CDIAL lists under Sanskrit \emph{piṃṣati}, \textquotedblleft{}grinds\textquotedblright{} (CDIAL 8142): to have someone grind.
\end{cousinweb}

\begin{grammarlens}[title={What you've built}]
One more word for describing what you lived and felt.
\end{grammarlens}

\subsection*{Your turn}
\begin{itemize}
\raggedright
  \item \emph{Say it:} \emph{pisvānā}
  \item \emph{Say it:} \emph{pisvānā}, once more
  \item \emph{Say it:} \emph{ammī ne gandum kā āṭā pisvāyā.} — Mum had the wheat ground into flour.
  \item \emph{Recall:} say the Urdu for to pound, then the Urdu for a mortar (for pounding), then say \emph{pisvānā} again
\end{itemize}

\begin{tcolorbox}[breakable,colback=teal!4,colframe=teal!35!black,title={Before you move on}]
What does \ur{پسوانا} mean? (\textquotedblleft{}To have (something) ground\textquotedblright{}.) Say it once more.
\end{tcolorbox}
\section[\texorpdfstring{\ur{نچوڑنا} (nichoṛnā)}{نچوڑنا (nichoṛnā)}]{\ur{نچوڑنا} (nichoṛnā) — to squeeze out (juice or water)}

\label{lesson:UR-C510-nichorna}



\begin{quote}
Before the new one: say the Urdu for a mortar (for pounding), then the Urdu for to have (something) ground.
\end{quote}

\subsection*{You'll want to know: \ur{نچوڑنا}}
\textbf{\ur{نچوڑنا}} — \emph{nichoṛnā} — \textquotedblleft{}to squeeze out (juice or water)\textquotedblright{}.

\begin{quote}
\ur{اس} \ur{نے} \ur{سلاد} \ur{پر} \ur{لیموں} \ur{نچوڑا۔}
\end{quote}

\begin{quote}
\emph{us ne salād par lemū̃ nichoṛā.}
\end{quote}

\begin{quote}
She squeezed a lemon over the salad.
\end{quote}

Past, with \ur{نے} (he or she …): \ur{اس} \ur{نے} \ur{تولیہ} \ur{نچوڑا}, \emph{us ne tauliyā nichoṛā}.

\begin{cousinweb}
\textbf{\ur{نچوڑنا}} (\emph{nichoṛnā}) is inherited: CDIAL lists the Hindi cousin \emph{nicoṛnā}, \textquotedblleft{}to squeeze out\textquotedblright{}, under \emph{*niscōṭayati}, a form it reconstructs (CDIAL 7449).
\end{cousinweb}

\begin{grammarlens}[title={What you've built}]
One more word for describing what you lived and felt.
\end{grammarlens}

\subsection*{Your turn}
\begin{itemize}
\raggedright
  \item \emph{Say it:} \emph{nichoṛnā}
  \item \emph{Say it:} \emph{nichoṛnā}, once more
  \item \emph{Say it:} \emph{us ne salād par lemū̃ nichoṛā.} — She squeezed a lemon over the salad.
  \item \emph{Recall:} say the Urdu for a mortar (for pounding), then the Urdu for to have (something) ground, then say \emph{nichoṛnā} again
\end{itemize}

\begin{tcolorbox}[breakable,colback=teal!4,colframe=teal!35!black,title={Before you move on}]
What does \ur{نچوڑنا} mean? (\textquotedblleft{}To squeeze out (juice or water)\textquotedblright{}.) Say it once more.
\end{tcolorbox}
